% =========================================================
\chapter{Manual de usuario}
\label{ap:manual_usuario}
% =========================================================

El sistema descrito en el \cref{cap:resultados} lo opera en el
laboratorio una persona que no programa. Para ella se redactó un manual
de usuario completo, que se entrega como documento aparte
(\texttt{documentos/manual\_usuario/manual\_usuario.pdf}) y corresponde
a la versión congelada del software: rama \texttt{main}, commit
\texttt{c67ae06}. Este apéndice resume su contenido, en particular el
flujo de operación, los diálogos que el operador encontrará y las
limitaciones conocidas, contrastados con el código de esa versión; los
detalles de cada pantalla y las hojas de registro están en el manual
completo.

\section*{Flujo de operación}

\begin{enumerate}[leftmargin=*, itemsep=0pt, topsep=2pt]
  \item \textbf{Verificación previa.} Osciloscopio encendido, con su
    servidor de red iniciado y el cable Ethernet conectado; láser y
    ESP32 conectados por USB, con los sensores S1--S4 en la muestra.
    Escalas, posición horizontal y disparo se ajustan en el propio
    osciloscopio: la aplicación no los modifica, solo los lee y los
    registra con cada medición.
  \item \textbf{Arranque.} Se ejecuta \texttt{iniciar.bat}, que abre
    la aplicación junto con una consola que no debe cerrarse, y en la
    pantalla de inicio se elige \emph{Láser + Osciloscopio}, la única
    ventana con temperatura y secuencias automáticas.
  \item \textbf{Conexión.} La primera vez, o si cambia un puerto, el
    diálogo \emph{Conexión} fija los puertos COM del ESP32 y del láser
    y la IP del osciloscopio, y los guarda. Cada equipo tiene un
    indicador: verde (conectado), amarillo (reconectando), rojo (sin
    conexión) o gris (sin estado todavía).
  \item \textbf{Geometría.} La separación hidrófono--haz y el volumen
    vertido se capturan antes de la primera medición: se registran al
    crearse la sesión y en ese momento los campos se bloquean.
  \item \textbf{Láser.} Se fijan sus parámetros, se pulsa
    \emph{Aplicar parámetros} y después \emph{Iniciar}, que pide
    confirmar que las protecciones están en su lugar. Con el láser en
    marcha los parámetros quedan bloqueados.
  \item \textbf{Osciloscopio.} Se elige el canal (CH1 para el
    hidrófono) y el modo \emph{Sample} o \emph{Average} con su número
    de promedios, y se pulsa \emph{Aplicar parámetros}.
  \item \textbf{Captura manual.} \emph{Capturar} muestra la señal para
    revisarla y \emph{Guardar} la registra. El primer guardado pide la
    carpeta base y crea en ella la sesión \texttt{AAAAMMDD\_HHMMSS};
    la secuencia automática solo se habilita con una sesión activa,
    que se obtiene al guardar una captura manual o al abrir un CSV
    existente.
  \item \textbf{Secuencia automática.} \emph{Por tiempo}: número de
    mediciones e intervalo, de al menos \SI{15}{\second}.
    \emph{Por temperatura}: temperatura inicial mayor que la final y
    paso mayor que cero, sobre una muestra que se enfría sola. Durante
    la secuencia el panel del osciloscopio queda bloqueado y la
    aplicación impone el modo \emph{Average} (100 promedios por tiempo;
    \num{10000} por temperatura), que restaura al terminar. Al
    concluir, por término normal, detención o aborto, el láser recibe
    las órdenes de modo seguro.
  \item \textbf{Revisión de datos.} En \emph{Visualizar datos} las
    filas con \texttt{error\_flag = 1} aparecen resaltadas; su causa
    está en \texttt{error\_desc}. Desde ahí se reabre una sesión
    (\emph{Abrir CSV\ldots}), se exporta o se convierte a MATLAB.
  \item \textbf{Cierre.} Se detiene el láser y se sale de la ventana; al
    cerrarla, la aplicación envía el modo seguro al láser si sigue
    conectado. La carpeta de la sesión se copia completa.
\end{enumerate}

\section*{Limitaciones conocidas}

\begin{itemize}[leftmargin=*, itemsep=0pt, topsep=2pt]
  \item \textbf{Inactividad.} El diálogo «¿Sigues usando el láser?» es un
    recordatorio y no sustituye la supervisión del operador
    (\cref{sec:protecciones}).
  \item \textbf{Modo seguro sin confirmación.} Las órdenes de modo
    seguro se envían sin esperar respuesta del láser; el operador debe
    comprobar visualmente que quedó detenido. El diálogo de dispositivo crítico
    desconectado no se ha verificado en hardware.
  \item \textbf{Modo de adquisición inicial.} Al abrir la ventana
    aparece seleccionado \emph{Sample} sin haberse leído del
    osciloscopio; aplicar sin revisarlo puede cambiar el modo del
    instrumento.
  \item \textbf{Metadatos de sesión.} La geometría y los parámetros del
    láser se escriben en el archivo de metadatos una sola vez, al crear
    la sesión; los parámetros vigentes en cada medición son los de las
    columnas del CSV.
  \item \textbf{Ajustes del osciloscopio.} Mover escalas o disparo
    durante una secuencia no la interrumpe: el cambio queda registrado
    fila por fila, puede marcar con \texttt{error\_flag} la captura en
    curso y mezcla condiciones dentro de la sesión.
\end{itemize}

\section*{Diálogos principales}

\noindent\begin{minipage}{\textwidth}
  \centering
  \footnotesize
  \captionof{table}{Diálogos que el operador encuentra con mayor
    frecuencia y respuesta recomendada.}
  \label{tab:dialogos_manual}
  \renewcommand{\arraystretch}{1.05}
  \setstretch{1.0}
  \noindent\begin{tabularx}{\textwidth}{>{\raggedright\arraybackslash}p{0.22\textwidth}
      >{\raggedright\arraybackslash}p{0.30\textwidth}
      >{\raggedright\arraybackslash}X}
    \toprule
    \textbf{Diálogo} & \textbf{Cuándo aparece} & \textbf{Qué responder}\\
    \midrule
    Dispositivo sin conexión & Al abrir la ventana, si algún equipo no
      responde & Revisar cable y puerto o IP en \emph{Conexión}.\\
    Encender el láser & Al pulsar \emph{Iniciar} & Aceptar solo con las
      protecciones colocadas.\\
    Geometría incompleta & Al capturar sin sesión creada y con campos
      vacíos & \emph{Cancelar y llenarlos} (opción por omisión) o
      \emph{Continuar sin ellos}.\\
    Error de captura / Datos inválidos & La captura no devolvió una
      señal utilizable & Revisar indicador del osciloscopio, disparo y
      láser; repetir.\\
    Guardar en sesión reabierta & Primer guardado o secuencia tras
      \emph{Abrir CSV\ldots} & \emph{Yes} solo para continuar esa
      sesión; la opción por omisión es \emph{No}.\\
    Configuración no realizable & Paso de temperatura nulo, o intervalo
      insuficiente para promediar y transferir & Corregir el parámetro
      indicado.\\
    Secuencia completada / detenida / abortada & Al terminar la
      secuencia & Revisar las filas marcadas o el motivo del aborto.\\
    ¿Sigues usando el láser? & Un minuto sin actividad, con el láser en
      marcha y sin secuencia & \emph{Yes} para continuar; \emph{No}
      activa el modo seguro.\\
    Dispositivo crítico desconectado & El láser o el osciloscopio
      no reconecta tras tres intentos & Leer la lista: «enviado» indica que la orden salió,
      no que el láser la confirmó.\\
    Reinicie la aplicación / Reinicio necesario & Un stop de emergencia
      sin confirmación del fin de la secuencia & Cerrar con la X y
      volver a abrir.\\
    Láser encendido & Al pulsar \emph{Volver} con el láser en marcha &
      \emph{Cancel} permanece en la ventana; con \emph{Yes} o \emph{No}
      se sale y se envía el modo seguro.\\
    \bottomrule
  \end{tabularx}
\end{minipage}
